\section*{PHÂN CÔNG CÔNG VIỆC}
\addcontentsline{toc}{section}{PHÂN CÔNG CÔNG VIỆC}

\begin{table}[H]
\centering
\captionsetup{justification=centering}
\caption{Bảng phân công công việc của các thành viên}
\label{tab:phancong}
\begingroup
\small
\setstretch{1.08}
\renewcommand{\arraystretch}{1.0}
\setlength{\tabcolsep}{5pt}
\setlength{\extrarowheight}{2pt}
\begin{tabularx}{\textwidth}{|c|>{\raggedright\arraybackslash}p{2.4cm}|>{\centering\arraybackslash}p{1.8cm}|>{\raggedright\arraybackslash}X|}
\hline
\multicolumn{1}{|c|}{\textbf{STT}} &
\multicolumn{1}{c|}{\textbf{Họ và tên}} &
\multicolumn{1}{>{\centering\arraybackslash}m{1.8cm}|}{\textbf{Mã sinh viên}} &
\multicolumn{1}{c|}{\textbf{Nhiệm vụ được giao}} \\
\hline
1 & Nguyễn Quang Anh & 24001649 &
Trưởng nhóm\newline
- Lập kế hoạch, phân chia công việc, điều phối tiến độ và tổng hợp, rà soát toàn bộ báo cáo.\newline
- Nghiên cứu bài báo gốc; phụ trách viết Phần 1--3 (giới thiệu, cơ sở lý thuyết năm mô hình song song, phương pháp đánh giá và thiết kế prompt).\newline
- Xây dựng bộ demo Mandelbrot (\code{Demo/})\ifdefined\isExtendedReport\ và kiểm chứng bổ sung trên Gemini 4 Hydrogen (Mục~\ref{subsec:gemini4})\fi.\newline
- Hoàn thiện định dạng LaTeX (trang bìa, mục lục, phụ lục, tài liệu tham khảo) và thuyết trình Phần 2: Kiến trúc thực nghiệm (kèm minh họa demo). \\[3pt]
\hline
2 & Nguyễn Như Bảo & 24001655 &
Thành viên\newline
- Nghiên cứu bài báo gốc; tổng hợp, đối chiếu số liệu và phụ trách viết Phần 4 (tỉ lệ thành công, 19 lỗi, phân tích LOC/COCOMO và khả năng mở rộng).\newline
- Phụ trách viết Phần 5--6 (thảo luận, đánh giá hiệu năng, bài học thực tiễn và kết luận, hướng phát triển).\newline
- Chuẩn hóa hệ thống bảng biểu, hình minh họa, danh mục từ viết tắt và rà soát độ chính xác trích dẫn tài liệu tham khảo.\newline
- Phối hợp cùng trưởng nhóm hiệu đính bản báo cáo cuối cùng và thuyết trình Phần 4: Hiệu năng \& tổng kết. \\[3pt]
\hline
3 & Đồng Đức Dương & 24001659 &
Thành viên\newline
- Nghiên cứu nội dung báo cáo, tập trung vào kết quả biên dịch, thực thi và phân loại các nhóm lỗi C++/OpenMP/MPI.\newline
- Thiết kế slide trình bày kết quả đánh giá các LLM, hệ thống phân loại lỗi thực tế và minh họa các đoạn mã lỗi tiêu biểu.\newline
- Biên tập bảng tổng hợp, hình ảnh trực quan và phối hợp cùng Trần Lý Dũng đồng bộ phong cách cho toàn bộ bộ slide.\newline
- Xây dựng kịch bản, luyện tập và thuyết trình Phần 3: Phân tích lỗi thực tế. \\[3pt]
\hline
4 & Trần Lý Dũng & 24001658 &
Thành viên\newline
- Nghiên cứu nội dung báo cáo, tập trung vào đặt vấn đề, bài toán Mandelbrot, các cơ chế song song và tổng kết hiệu năng.\newline
- Thiết kế slide cho phần giới thiệu bài toán, kiến trúc thực nghiệm và phần hiệu năng, tổng kết, khuyến nghị.\newline
- Biên tập sơ đồ quy trình, biểu đồ mở rộng và tối ưu bố cục trình chiếu cho bộ slide chung của nhóm.\newline
- Xây dựng kịch bản, luyện tập và thuyết trình Phần 1: Giới thiệu bài toán. \\[3pt]
\hline
\end{tabularx}
\endgroup
\end{table}
